\documentclass[../../../include/open-logic-section]{subfiles}

\begin{document}

\olfileid{sth}{ord-arithmetic}{add}
\olsection{ক্রমসংখ্যার যোগ}

ধরি, আমরা $\alpha$ ও $\beta$ যোগ করতে চাই।
$\alpha$-র একটি {অনুলিপির} ঠিক পরে $\beta$-র একটি
অনুলিপি বসালেই হয়। (অনুলিপি নিতে হয়, কারণ
\olref[ordinals][basic]{ordinalsaresubsets} থেকে জানি,
$\alpha \subseteq \beta$ অথবা $\beta \subseteq \alpha$।)
স্বজ্ঞাগতভাবে এতে প্রথমে $\alpha$-দীর্ঘ একটি ধাপক্রম,
\emph{তার পরে} $\beta$-দীর্ঘ আরেকটি ধাপক্রম পাওয়া যায়।
ফলের ক্রমটি সুক্রমবিন্যস্ত; তাই
\olref[ordinals][ordtype]{thmOrdinalRepresentation}
অনুযায়ী এটি একটি (অনন্য) ক্রমসংখ্যার সঙ্গে সমরূপ।
সেই ক্রমসংখ্যাকে $\alpha$ ও $\beta$-র \emph{যোগফল}
ধরা যায়।

ক্রমসংখ্যার যোগের পেছনে এই স্বজ্ঞাই কাজ করে।
কঠোর সংজ্ঞা দিতে প্রথমে সেটের \emph{অনুলিপি} নেওয়ার
ধারণা ব্যবহার করি। কোন বস্তুটি কোথা থেকে এসেছে তা
চিহ্নিত রাখতে এখানে যথেচ্ছ চিহ্ন $0$ ও $1$ নিই:

\begin{defn}\ollabel{defdissum}
$A$ ও $B$-র \emph{বিযুক্ত যোগফল} হলো $A \disjointsum B =
(A\times \{0\}) \cup (B \times \{1\})$।
\end{defn}

এবার ক্রমসংখ্যার জোড়গুলির উপর একটি ক্রম সংজ্ঞায়িত করি:

\begin{defn}
যেকোনো ক্রমসংখ্যা $\alpha_1, \alpha_2, \beta_1, \beta_2$-এর জন্য ধরি:
\begin{align*}
	\tuple{\alpha_1, \alpha_2} \rlexless \tuple{\beta_1, \beta_2}\text{ যদি এবং কেবল যদি }&
	\text{হয় $\alpha_2 \in \beta_2$}\\
	& \text{নয়তো $\alpha_2 = \beta_2$ এবং $\alpha_1 \in \beta_1$}
\end{align*}
\end{defn}

এটি \emph{বিপরীত অভিধানিক} ক্রম: আগে দ্বিতীয়
স্থানাঙ্ক, তার পরে প্রথম স্থানাঙ্ক দিয়ে তুলনা করা হয়।
আমাদের উদ্দেশ্য ছিল, $\alpha$-র একটি অনুলিপির পরে
$\beta$-র একটি অনুলিপি বসিয়ে যে ক্রম পাওয়া যায়,
তার ক্রমপ্রকারকে $\alpha \ordplus \beta$ বলা। তাই ধরি:

\begin{defn}\ollabel{defordplus}
যেকোনো ক্রমসংখ্যা $\alpha$, $\beta$-র যোগফল
$\alpha \ordplus \beta = \ordtype{\alpha \disjointsum \beta,
\rlexless}$।
\end{defn}
\noindent
এখানে আমরা সংকেতের সামান্য শিথিল ব্যবহার করেছি;
% BN-SRC-801: সম্পর্কটি ক্ষেত্রের কার্তেসীয় বর্গে; চিহ্নিত অনুলিপির সদস্যকে সম্পর্কের যুগল ভাবা নয়।
কঠোরভাবে ``$\rlexless$''-এর বদলে
``$\Setabs{\tuple{x,y}\in (\alpha\disjointsum\beta)\times(\alpha\disjointsum\beta)}{x \rlexless y}$''
লেখা উচিত। সংক্ষিপ্ততার জন্য পরেও এভাবে সংকেত ব্যবহার করব।

নিচের ফলটি ও
\olref[ordinals][ordtype]{thmOrdinalRepresentation}
মিলে আমাদের সংজ্ঞাটি সুগঠিত বলে নিশ্চিত করে:

\begin{lem}\ollabel{ordsumlessiswo}
যেকোনো ক্রমসংখ্যা $\alpha$ ও $\beta$-র জন্য
$\tuple{\alpha \disjointsum \beta, \rlexless}$
একটি সুক্রমবিন্যাস।
\end{lem}

\begin{proof}
$\alpha \disjointsum \beta$-তে $\rlexless$-এর অধীনে
যেকোনো দুটি ভিন্ন উপাদান তুলনাযোগ্য। এটি সুভিত্তিক
দেখাতে একটি অশূন্য $X \subseteq \alpha \disjointsum \beta$
নিই। $X$-এর যে উপসেটে দ্বিতীয় স্থানাঙ্কটি সম্ভাব্য
ক্ষুদ্রতম, সেটি $Y$ হোক; অর্থাৎ
$Y = \Setabs{\tuple{\gamma, i} \in X}{(\forall
\tuple{\delta, j} \in X)i \leq j}$। এবার $Y$ থেকে
ক্ষুদ্রতম প্রথম স্থানাঙ্কের উপাদানটি নিই।
\end{proof}
\noindent
অতএব ক্রমসংখ্যার যোগের একটি সুস্পষ্ট, সরাসরি সংজ্ঞা
পেলাম। এবার একটি প্রত্যাশিত ফল (\olref[z][infinity-again]{defnomega}
অনুযায়ী $1 = \{0\}$):
\begin{prop}
যেকোনো ক্রমসংখ্যা $\alpha$-র জন্য
$\alpha \ordplus 1 = \ordsucc{\alpha}$।
\end{prop}

\begin{proof}
% BN-SRC-447: চিহ্নিত অনুলিপি দুটির সংযোগ; আবার বিযুক্ত যোগ নয়।
$\ordsucc{\alpha} = \alpha \cup \{\alpha\}$ থেকে
$\alpha\disjointsum1 = (\alpha \times \{0\})
\cup (\{0\} \times \{1\})$-তে সমরূপতা $f$ নিই,
যেখানে $\gamma \in \alpha$ হলে $f(\gamma) =
\tuple{\gamma, 0}$ এবং $f(\alpha) = \tuple{0, 1}$।
\end{proof}
\noindent
এ ছাড়া দেখানো সহজ, যোগ কয়েকটি পুনরাবৃত্তিমূলক
শর্ত মেনে চলে:

\begin{lem}\ollabel{ordadditionrecursion}
যেকোনো ক্রমসংখ্যা $\alpha, \beta$-র জন্য:
\begin{align*}
	\alpha\ordplus 0 &= \alpha\\
	\alpha \ordplus (\beta\ordplus 1) &= (\alpha \ordplus \beta) \ordplus 1\\
	\alpha \ordplus \beta &= \supstrict_{\delta < \beta}(\alpha \ordplus \delta)
	&& \text{যদি $\beta$ সীমা ক্রমসংখ্যা হয়}
\end{align*}
\end{lem}

\begin{proof}
প্রতিটি ক্ষেত্র আলাদা করে যাচাই করি। প্রথমে:
% BN-SRC-448: শূন্য কার্তেসীয় গুণফল খালি; শূন্যকে নিয়ে একপদী সেট নয়।
\begin{align*}
	\alpha \ordplus 0
	&= \ordtype{(\alpha \times \{0\}) \cup (0 \times \{1\}), \rlexless} \\
	&= \ordtype{(\alpha \times \{0\}) \cup \emptyset, \rlexless}\\
	&= \alpha\\
	\alpha \ordplus (\beta \ordplus 1)
	%&= \ordtype{(\alpha\times \{0\}) \cup (\ordtype{\beta \ordplus 1}\times \{1\}), \rlexless} \\
	&= \ordtype{(\alpha\times \{0\}) \cup (\ordsucc{\beta}\times \{1\}), \rlexless} \\
	&= \ordtype{(\alpha\times \{0\}) \cup (\beta \times \{1\}), \rlexless} \ordplus 1\\
	&= (\alpha \ordplus \beta) \ordplus 1
\end{align*}
এবার $\beta \neq \emptyset$ একটি সীমা ক্রমসংখ্যা
হোক। $\delta < \beta$ হলে $\delta\ordplus 1 < \beta$-ও
হয়; তাই $\alpha \ordplus \delta$ হলো $\alpha
\ordplus \beta$-র যথার্থ প্রারম্ভিক খণ্ড। অতএব
$\alpha \ordplus \beta$ হলো
$X = \Setabs{\alpha \ordplus \delta}{\delta < \beta}$-র
একটি যথার্থ ঊর্ধ্বসীমা। আবার $\alpha \leq \gamma <
\alpha \ordplus \beta$ হলে কোনো $\delta < \beta$-র
জন্য $\gamma = \alpha \ordplus \delta$। সুতরাং
$\alpha \ordplus \beta = \supstrict_{\delta < \beta}
(\alpha\ordplus \delta)$।
\end{proof}

কিন্তু এখানে একটি লক্ষণীয় ব্যাপার আছে। ক্রমসংখ্যার যোগ
সংজ্ঞায়িত করতে আমরা \emph{বিকল্পভাবে} শুধু
ট্রান্সফাইনাইট পুনরাবৃত্তি-উপপাদ্য ব্যবহার করতে পারতাম:
\olref{ordadditionrecursion}-এর শর্তগুলিকেই পুনরাবৃত্তি-সমীকরণ
হিসেবে দিতাম (যদিও ``$\beta \ordplus 1$''-এর বদলে
``$\ordsucc{\beta}$'' লিখতাম)।

অতএব ক্রমসংখ্যার উপর ক্রিয়া সংজ্ঞায়িত করার দুটি পথ আছে।
সুস্পষ্ট একটি সুক্রমবিন্যস্ত সেট নির্মাণ করে তার ক্রমপ্রকার
নিলে পাই \emph{সংশ্লেষমূলক} সংজ্ঞা। শুধু
পুনরাবৃত্তি-সমীকরণ দিলেও পাই \emph{পুনরাবৃত্তিমূলক}
সংজ্ঞা। ঠিকভাবে করলে দুটি পথেই একই ফল আসে। কারণ
\olref[ordinals][ordtype]{thmOrdinalRepresentation}
নিশ্চিত করে, এই পুনরাবৃত্তি-সমীকরণগুলি \emph{অনন্য}
ক্রমসংখ্যা নির্ধারণ করে।
% BN-SRC-802 TARGET-CORRECTION-BEGIN
\emph{উৎস-সংশোধন-নোট।} ক্রমপ্রকারের অনন্যতা এবং
পুনরাবৃত্তি-সমীকরণের অনন্যতার কারণ পৃথক। পরেরটিতে
একই $\alpha$ রেখে দুটি মান-পরিবার তুলনা করি। শূন্যে
সমান, উত্তরসূরিতে আগের সমান মানের একই উত্তরসূরি,
সীমায় আগের সমান পরিবারের একই ঊর্ধ্বসীমা। তাই
ট্রান্সফাইনাইট আবেশে প্রত্যেক সূচকে মান সমান।
সরাসরি ক্রমপ্রকার ওই সমীকরণ পূরণ করে বলেই দুই সংজ্ঞা মেলে।
% BN-SRC-802 TARGET-CORRECTION-END

অনেক দিক থেকে ক্রমসংখ্যার যোগ স্বাভাবিক সংখ্যার যোগের
মতোই। যেমন, আমরা প্রমাণ করতে পারি:

\begin{lem}\ollabel{ordinaladditionisnice}
$\alpha, \beta, \gamma$ ক্রমসংখ্যা হলে:
\begin{enumerate}
	\item\ollabel{ordaddition1} $\beta < \gamma$ হলে $\alpha
	\ordplus \beta < \alpha \ordplus \gamma$;
	\item\ollabel{ordaddition2} $\alpha \ordplus \beta =
	\alpha\ordplus \gamma$ হলে $\beta = \gamma$;
	\item\ollabel{ordaddition3} $\alpha \ordplus (\beta \ordplus
	\gamma) = (\alpha \ordplus \beta) \ordplus \gamma$;
	অর্থাৎ যোগ সম্বন্ধীয়: বন্ধনী বদলালেও ফল বদলায় না;
	\item\ollabel{ordaddition4} $\alpha \leq \beta$ হলে $\alpha
	\ordplus \gamma \leq \beta \ordplus \gamma$।
\end{enumerate}
\end{lem}

\begin{proof}
\olref{ordaddition3} প্রমাণ করি; বাকিগুলি অনুশীলনের জন্য
রাখছি। $\gamma$-র উপর সরল ট্রান্সফাইনাইট আবেশে,
\olref{ordadditionrecursion} ব্যবহার করে প্রমাণটি হবে।
$\gamma = 0$ হলে:
\[
(\alpha \ordplus \beta) \ordplus 0 = \alpha \ordplus \beta
= \alpha \ordplus (\beta \ordplus 0)
\]
$\gamma = \delta\ordplus 1$ হলে আবেশ-অনুমান ধরি:
$(\alpha \ordplus \beta) \ordplus \delta = \alpha
\ordplus (\beta \ordplus \delta)$। এবার
\olref{ordadditionrecursion} তিনবার প্রয়োগ করে পাই:
\begin{align*}
	(\alpha \ordplus \beta) \ordplus (\delta \ordplus 1)
	&= ((\alpha \ordplus \beta) \ordplus \delta)\ordplus 1\\
	&= (\alpha \ordplus (\beta \ordplus \delta)) \ordplus 1\\
	&= \alpha \ordplus ((\beta \ordplus \delta)\ordplus 1)\\
	&= \alpha \ordplus (\beta \ordplus (\delta\ordplus 1))
\end{align*}
$\gamma$ সীমা ক্রমসংখ্যা হলে আবেশ-অনুমান ধরি:
$\delta \in \gamma$ হলেই $(\alpha \ordplus \beta)
\ordplus \delta = \alpha \ordplus (\beta \ordplus
\delta)$। তখন:
\begin{align*}
	(\alpha \ordplus \beta) \ordplus \gamma
	&= \supstrict_{\delta < \gamma}((\alpha \ordplus \beta) \ordplus \delta) \\
	&= \supstrict_{\delta < \gamma}(\alpha \ordplus (\beta \ordplus \delta))\\
	&= \alpha \ordplus \supstrict_{\delta < \gamma}(\beta \ordplus \delta)\\
	&= \alpha \ordplus (\beta \ordplus \gamma)
\end{align*}
% BN-SRC-803 TARGET-CORRECTION-BEGIN
\emph{উৎস-সংশোধন-নোট।} সীমার তৃতীয় সমতায় অনুক্ত
সহশেষিতার যুক্তি আছে। $Y=\Setabs{\beta\ordplus\delta}{\delta<\gamma}$
অশূন্য এবং সর্বাধিক সদস্যহীন: $\delta\ordplus1<\gamma$
ও যোগের কঠোর বৃদ্ধিতে প্রতিটি সদস্যের পরে আরেকটি আছে।
তাই $\eta=\supstrict Y=\bigcup Y$ একটি সীমা।
প্রত্যেক $\xi<\eta$-র পরে বা সমান কোনও $y\in Y$ আছে;
ডান দিকে যোগের বৃদ্ধির ধর্মে $\alpha\ordplus\xi$
সেই $\alpha\ordplus y$-র নিচে বা সমান। ফলে সীমা-সমীকরণের
সব $\xi<\eta$-র ঊর্ধ্বসীমা শুধু সহশেষী $Y$ নিলেও বদলায় না।
শূন্য পরিবারে এ ধরনের বিনিময় নির্বিশেষে সত্য নয়।
% BN-SRC-803 TARGET-CORRECTION-END
\end{proof}

\begin{prob}
\olref[sth][ord-arithmetic][add]{ordinaladditionisnice}-এর
অবশিষ্ট দাবিগুলি প্রমাণ করুন।
\end{prob}

এই সব দিক থেকে ক্রমসংখ্যার যোগ বেশ পরিচিত মনে হবে।
কিন্তু একটি গুরুত্বপূর্ণ দিক থেকে তা স্বাভাবিক সংখ্যার যোগের
\emph{মতো নয়}।

\begin{prop}\ollabel{ordsumnotcommute}
ক্রমসংখ্যার যোগ {বিনিময়ী নয়};
$1 \ordplus \omega = \omega < \omega \ordplus 1$।
\end{prop}

\begin{proof}
লক্ষ করুন, $1 \ordplus \omega = \supstrict_{n < \omega}
(1 \ordplus n) = \omega \in \omega \cup \{\omega\} =
\ordsucc{\omega} = \omega \ordplus 1$।
\end{proof}
\noindent
প্রথমে এতে অবাক লাগতে পারে, কিন্তু \emph{লাগার কথা নয়}।
একদিকে, $1 \ordplus \omega$-তে সব স্বাভাবিক সংখ্যার
\emph{আগে} একটি বাড়তি উপাদান বসানো হচ্ছে।
\olref[sfr][infinite][hilbert]{sec}-এ হিলবার্টের হোটেল
নিয়ে যেভাবে যুক্তি করেছি, তাতে স্বজ্ঞাগতভাবে এই প্রথম
উপাদানটি সামগ্রিক ক্রমপ্রকার বদলায় না। অন্যদিকে,
$\omega \ordplus 1$-তে সব স্বাভাবিক সংখ্যার
\emph{পরে} একটি বাড়তি উপাদান বসানো হচ্ছে।
সেটি সম্পূর্ণ ভিন্ন ব্যাপার!

\end{document}
